% ============================================================
%  §9  Relation to the Note Series
% ============================================================

\section{Relation to the Note Series}
\label{sec:note}

The companion note series (L1 layer) developed the intuitive picture
of conceptual landscapes in accessible prose, using the vocabulary of
sedimentation, flow, and terrain.
This section makes the correspondence between the note-layer
descriptions and the formal structure of the present paper explicit.

\begin{center}
\renewcommand{\arraystretch}{1.5}
\begin{tabular}{ll}
\toprule
\textbf{Note-layer description (J: Japanese original)} &
\textbf{Formal correspondence} \\
\midrule
\textit{log no chindenˆ} ``log sedimentation''     & Deepening of $\Phi$ via $\sigma_{\mathrm{drive}}$ \\
\textit{anteiten} ``stable point''                  & Locally stable region: $\nabla\Phi\approx 0$ \\
\textit{gainen no kubomi} ``conceptual hollow''     & Potential well in $\Phi$ \\
\textit{zentai koubai} ``overall gradient''         & Information potential $\Phi = K*I$ \\
\textit{tooriyasusa} ``passability''                & $-\Phi(x,\tau)$ (high passability $=$ low $\Phi$) \\
\textit{tani} ``valley''                            & Deep well in $\Phi$ \\
\textit{kabe} ``wall''                              & Ridge region of high $\Phi$ \\
\textit{rikai shita kankaku} ``sense of understanding'' & Reduction of traversal cost $\int|\nabla\Phi|dx$ \\
\textit{shikou no kuse} ``habitual thought pattern'' & Limit-cycle trajectory of $\JI$ \\
\textit{GPU-teki na souchi} ``GPU-like device''     & Sub-threshold operation of $\Phi$ on $\JI$ \\
\textit{gengo ryuushi} ``language particle''        & Discrete token $\ell_i$ coupling to $\nabla\Phi$ \\
\textit{asshuku} ``compression''                    & Quantization operator $\mathcal{Q}$ \\
\textit{shakai-teki kyouyu chikei} ``shared social terrain'' & Social $\Phistar$ maintained across agents \\
\bottomrule
\end{tabular}
\end{center}

\subsection{What the Note Layer Does Not Claim}

The note series deliberately avoided technical formalism and did not
commit to specific functional forms.
Several descriptions in the note layer are, from the standpoint of
the present paper, \emph{interpretations} of field-theoretic quantities
rather than definitions.

In particular:
\begin{itemize}
  \item The note descriptions of ``meaning'' as an emergent property
        of repeated traversal correspond formally to the depth of
        potential wells, not to any semantic object.
  \item The note descriptions of ``understanding'' as a change in
        ease correspond to the traversal cost condition
        \eqref{eq:understanding}.
  \item The CPU/GPU distinction introduced in the Consciousness note
        corresponds to the fast-lane / slow-lane decomposition that
        will be formalized in the companion Consciousness Theory paper.
\end{itemize}

\subsection{The L3 → L1 Direction}

The present paper (L3) does not revise the note series (L1).
The relation is one of support, not replacement:
the formal structure established here provides the theoretical basis
that stabilizes the intuitive descriptions without overwriting them.

Where the note layer uses metaphor (terrain, flow, sedimentation),
the present paper provides the field-theoretic referents.
Where the note layer uses phenomenological description (understanding,
meaning, thought), the present paper provides structural conditions.
The layers are complementary; neither is reducible to the other.
